\scanpage{121}
\mspage{59}
% L01
\[
\boxed{\sigma^{*}(F_D)\xrightarrow{\varphi}\rho^{*}(F_U)}
\]
% L02
Now, \deleted{\unsure{this datum if $F_U$ is a}} \inserted{if we restrict ourselves to the} $F_U$ such that $F_U$ is
% L03
itself essentially$^{t}$ locally$^{t}$ constant, then the data
% L04
$F_U,F_D,\varphi$ involve only
% L05
essentially$^{t}$ locally constant sheaves, i.e.
% L06
fundamental topoi (multi-Galois)
% L07
associated with the schemes under consideration, \deleted{\unsure{namely}} thus
% L08
are described entirely in terms of the
% L09
diagrams
\diagram{\begin{tikzcd}[ampersand replacement=\&]
 \& B_{\pi_1,D^{*}} \arrow[dl,"{\deleted{\illegible{1}}\,\rho}"'] \arrow[dr,no head,"\sigma"] \& \\
 B_{\pi_1,U} \& \& B_{\pi_1,D}
\end{tikzcd}}
\margin{(*)}
% L10
-- moreover the case where $\varphi$ is an \uline{isomorphism} corresponds
% L11
precisely to the case of locally essentially
% L12
constant sheaves on $X$, i.e. of sheaves on
% L13
$B_{\pi_1,X}$ -- which fundamental topos
% L14
thus appears as the amalgamated
% L15
sum of the preceding diagram of topoi,
% L16
fitting into the square
\diagram{\begin{tikzcd}[ampersand replacement=\&]
 \& B_{\pi_1,D^{*}} \arrow[dl] \arrow[dr] \& \\
 B_{\pi_1,U} \arrow[dr] \& \& B_{\pi_1,D} \arrow[dl] \\
 \& B_{\pi_1,X} \&
\end{tikzcd}}

\scanpage{122}
% L01
Suppose for simplicity $U$ connected,
% L02
let us choose a universal covering \inserted{$\widetilde U$} of $U$, and a
% L03
universal covering \inserted{$\widetilde D_i^{*}$} of each $D_i^{*}$ -- whence a
% L04
universal covering \inserted{$\widetilde D_i$} of $D_i$ -- and isomorphisms
% L05
(= ``classes of paths'') between $\rho_{!}(\widetilde D_i^{*})$ and $\widetilde U$
% L06
-- thus the diagram of topoi \unsure{(*)},
% L07
-- or of fundamental groupoids --
% L08
is made explicit in terms of a diagram
% L09
\diagram{\begin{tikzcd}[ampersand replacement=\&,column sep=large]
 \& {\begin{gathered}
 (\operatorname{Gal}(\overline K_i/K_i))\\[-2pt]
 \Vert\\[-2pt]
 \mathcal E_i
 \end{gathered}} \arrow[dl,"\unsure{\rho_i|\varphi_i}"'] \arrow[dr,"\sigma_i"] \& \\
 {\begin{gathered}
 \mathcal E\\[-2pt]
 \Vert\\[-2pt]
 (\pi_1(U,\widetilde U)=\operatorname{Aut}_U(\widetilde U))
 \end{gathered}} \& \&
 {\begin{gathered}
 \Gamma_i\;\illegible{1}\\[-2pt]
 \Vert\\[-2pt]
 (\operatorname{Gal}(\overline{k(s)}/k(s)))
 \end{gathered}}
\end{tikzcd}}
% L10
and the datum of an $F$ as considered on
% L11
$X$ amounts to the datum of a system
% L12
$(E_U,(E_i)_{i\in I},(\varphi_i)_{i\in I})$, where $E_U$ is an $\mathcal E$-set,
% L13
$E_i$ a $\Gamma_i$-set, $(\forall i\in I)$, and $\varphi_i$ an
% L14
$\mathcal E_i$-homomorphism from $E_i$ to $E_U$ (i.e. a
% L15
a $\Gamma_i$-homomorphism
\[
\varphi_i:\ E_i\longrightarrow E_U^{\Pi_i},
\]

\scanpage{123}
\mspage{60}
% L01
where \deleted{\unsure{$\Pi_i$ is}} $\Pi_i=\unsure{\operatorname{Ker}(\sigma_i)}$ fits into the exact sequence
% L02
\[
1\longrightarrow\Pi_i\longrightarrow\mathcal E_i\longrightarrow\Gamma_i\longrightarrow1
\]

% L03
The case where the $E_i\longrightarrow E_U$ are isomorphisms,
% L04
i.e. the $E_i$ \unsure{are identified|being identified} \deleted{\illegible{1}} all with $E_U$,
% L05
with trivial action \unsure{of|of the} $\Pi_i$ on $E_U$, corresponds
% L06
to the case where $F$ is essentially locally \unsure{constant|constant} on $X$,
% L07
whence at once
% L08
\[
\pi_1(X)=\mathcal E=\pi_1(U)\ \big/\
\substack{\text{\unsure{normal} subgroup generated}\\
\text{by the }\rho_i(\Pi_i)}.
\]

% L09
Our aim is that of a Galois-theoretic calculation
% L10
(if possible) of the cohomology of $X$ for
% L11
the $F$ under consideration.
% L12
Let $B_{X,U}$ be the topos \inserted{whose sheaves are} triples
% L13
$(F_U,F_D,\varphi)$ as above, which thus maps
% L14
to the topos $B_{\pi_1,X}$, corresponding
% L15
to the pairs for which $\varphi$ is an isomorphism,
% L16
and receives the topos $X_{\mathrm{ét}}$ :

\scanpage{124}
% L01
\diagram{\begin{tikzcd}[ampersand replacement=\&]
X_{\mathrm{ét}} \arrow[r,"f"] \& B_{X,U} \arrow[r,"g"] \& B_{\pi_1,X}
\end{tikzcd}}
% L02
We ask ourselves the question

% L03
a) ``Explicit'' calculation, in terms of the cohomology
% L04
of profinite groups, of cohomology
% L05
in the abracadabra topos $B_{X,U}$. (N.B.
% L06
\deleted{\unsure{b) the calculation of}} the cohomology \deleted{\illegible{1}} of
% L07
$B_{\pi_1,X}$ is none other than the profinite
% L08
Galois cohomology of $\pi_1(X)=\mathcal E/\ldots$
% L09
calculated above\ldots).

% L10
b) Check whether for a torsion
% L11
sheaf $F$ on $B_{X,U}$ the canonical homomorphism
% L12
\[
H^{*}(B_{X,U},F)\longrightarrow
H^{*}(X_{\mathrm{ét}},f^{*}F)
\]
% L13
is an isomorphism.

% L14
N.B. Up to now, the hypothesis
% L15
$\dim X=1$ has not been used, nor the noetherian hypothesis
% L16
-- only the fact that $S$ is a closed discrete
% L17
subset, $X$ connected\ldots. For the calculations which

\scanpage{125}
\mspage{61}
% L01
follow, corresponding to the \deleted{\illegible{1}} case where the
% L02
$\unsure{D_s|O_s}$ are \unsure{gauges} with residue field $k(s)$
% L03
of characteristic $0$, we are going to suppose that
% L04
\[
\Pi_i\simeq\widehat{\mathbb Z}\qquad\text{(not canonically)}
\]
% L05
-- in fact, by Kummer theory we
% L06
have a system of Tate local systems
% L07
$\unsure{T_{\mathcal E}}$ \inserted{(on $\mathcal E$)}, $T_i$ \inserted{(on $\Gamma_i$)}, and \inserted{$\mathcal E_i$-}isomorphisms $T_{\mathcal E}\simeq T_i$, \deleted{and \illegible{1}}
% L08
(i.e. a $T_X$ on which $\unsure{\widetilde{\mathcal E}|\overline{\mathcal E}}=\mathcal E/\ldots$ operates)
% L09
and we shall have a canonical Kummer isomorphism
% L10
\[
T_i\xrightarrow[\sim]{\unsure{\kappa_i|k_i}}\Pi_i.
\]

% L11
We would like to be able to paraphrase,
% L12
on $B_{X,U}$, the well-known exact
% L13
cohomology sequence
% L14
\[
\longrightarrow H^i(X,F)\longrightarrow H^i(U,F)
\longrightarrow H^{i+1}_S(X,F)\longrightarrow\cdots
\]
% L15
relative to the open $U$ of the topos $X$, and its
% L16
closed ``complement'' $S$. Now, \deleted{\illegible{2}} \inserted{just as}
% L17
$X$ fits into a diagram of topoi
% L18
\diagram{\begin{tikzcd}[ampersand replacement=\&]
 \& D^{*}_{\mathrm{ét}} \arrow[dl] \arrow[dr] \& \& \\
 U_{\mathrm{ét}} \arrow[dr,no head] \& \& D_{\mathrm{ét}} \arrow[dl] \& S_{\mathrm{ét}} \arrow[l,"\unsure{\xi}"'] \\
 \& X_{\mathrm{ét}} \& \&
\end{tikzcd}},

\scanpage{126}
% L01
likewise $B_{X,U}$ fits into
% L02
\diagram{\begin{tikzcd}[ampersand replacement=\&]
 \& B_{\pi_1,D^{*}} \arrow[dl] \arrow[dr] \& \& \\
 B_{\pi_1,U} \arrow[dr] \& \& B_{\pi_1,D} \arrow[dl] \& B_{\pi_1,S} \arrow[l,"\approx","\unsure{\xi_{\mathrm{ét}}}"'] \\
 \& B_{X,U} \& \&
\end{tikzcd}}

% L03
[\uline{N.B.} the finite $\varprojlim$ and arbitrary $\varinjlim$
% L04
in $B_{X,U}$ i.e. for the $(F_U,F_D,\varphi)$, are
% L05
calculated ``term by term''] I say
% L06
that $B_{\pi_1,U}\longrightarrow B_{X,U}$ is identified with an
% L07
induction morphism, relative to the
% L08
object \deleted{\illegible{1}} \unsure{(again denoted} $U$ by abuse of notation)
% L09
of $B_{X,U}$ defined by
% L10
$F_U={}$final sheaf $e_U$, $F_D={}$initial sheaf $\varnothing_D$
% L11
(and $\varphi$ \unsure{is|being} then fixed !). Indeed, the
% L12
objects \inserted{of $B_{X,U}$} over $U$ are identified with the
% L13
$(F_U,F_D,\varphi)$ with $F_D=\varnothing_D$ (which $\varphi$ fixes !) so
% L14
they form a category equivalent to that of the
% L15
$F_U$, i.e. to $B_{\pi_1,U}$. We see that the residual
% L16
topos (being identified \deleted{with the} \inserted{with the category of} $F\deleted{_{\unsure{U}},\unsure{E|e}_{\unsure{i}}}=(F_U,F_D,\varphi)$ such

\scanpage{127}
\mspage{62}
% L01
that $F_U=e_U$ final sheaf, is likewise identified with the
% L02
category \inserted{$B_{\pi_1,D}$} of the $F_D$ \deleted{\illegible{2}} -- the
% L03
canonical functor ``inverse image on $B_{\pi_1,D}$ of the direct
% L04
image of $F$ on $B_{\pi_1,U}$'' being none other
% L05
$\sigma_*\rho^*$, so that we recover the description
% L06
\deleted{\illegible{1}} typical of Artin of a topos obtained
% L07
by ``gluing'' an \deleted{\unsure{object}} open and the
% L08
complementary closed.
% L09
We thus find an exact sequence
% L10
\[
\longrightarrow H^i(B_{X,U},F)\longrightarrow
H^i(B_{\pi_1,U},F_U)\longrightarrow
H^{i+1}_{S(\text{or }D)}(B_{X,U},F)\longrightarrow\cdots
\]
% L11
mapping into the analogous exact sequence
% L12
relative to $X_{\mathrm{ét}},U_{\mathrm{ét}},S_{\mathrm{ét}}$. \deleted{\unsure{For} \illegible{2}} \inserted{(We have}
% L13
still not used any special hypothesis on $S$\ldots).

% L14
To check that we have isomorphisms
% L15
at the level of the $H^i(B_{X,U},F)\longrightarrow H^i(X_{\mathrm{ét}},F)$,
% L16
it suffices by the five lemma to
% L17
prove it at the level of the
% L18
$H^i(B_{\pi_1,U},F)\longrightarrow H^i(U_{\mathrm{ét}},F)$ [i.e. \unsure{check|let us check} that the
% L19
cohomology of $U_{\mathrm{ét}}$ ``is calculated Galois-theoretically'']
% L20
\uline{and} at the level of the $H^i_S$. \deleted{\unsure{Now}} Let us look,

\scanpage{128}
% L01
first at the latter -- we have
% L02
$H^0_S(B_{X,U},E)={}$ set of sections of $F$ [i.e. of
% L03
pairs of a section of $F_U$ and
% L04
one of $F_D$, corresponding to each other
% L05
under $\varphi$, i.e. of an element \inserted{$x$} of
% L06
$E_U$ invariant under $\mathcal E$, and of
% L07
$x_i\in E_i$ invariant under the $\Gamma_i$,
% L08
such that $x=\varphi_i(x_i)\quad\forall i$]
% L09
such that $x=0$ \deleted{\unsure{of the|in} \unsure{systems}}
% L10
\[
=\prod_i(\operatorname{Ker}(E_i\longrightarrow E_U)
\]
% L11
we observe that this functor factors
% L12
through the functor $H^0_S$ relative to the ``inclusion''
% L13
\deleted{\unsure{$D^{*}\subset D$} \unsure{or} $B_{\pi_1,D^{*}}\longrightarrow B_{\pi_1,D}$}
% L14
analogous \deleted{\illegible{1} $B_{\pi_1,U}$ \illegible{2}}
% L15
\[
B_{\underbrace{\pi_1,\unsure{D^{*}|D_{*}}}_{\text{replacing }\pi_1,U}}\longrightarrow B_{D,D^{*}}\quad\ldots
\]
% L16
and the $H^i_S$ on $D_{X,U}$ are calculated like
% L17
\deleted{\unsure{which is none other than}}
% L18
those in the local situation $B_{D\deleted{\illegible{1}},D^{*}\deleted{\illegible{1}}}$ --
% L19
where we find the usual calculation in
% L20
terms of the group homomorphism $\mathcal E_i\longrightarrow\Gamma_i$.
\margin{which I have forgotten a bit !}

% L21
But in the present case, the hypothesis $\dim_s X=1$
% L22
at the points of $s\in S$, \deleted{and $X$ \unsure{regular}} implies that
% L23
the localized situation \inserted{of topoi} $D^{*}_{\unsure{a|\mathrm{ét}}}\hookrightarrow D_{\unsure{a|\mathrm{ét}}}$ is already entirely
% L24
defined in terms of the fundamental topoi \unsure{of the}

\scanpage{129}
\mspage{63}
% L01
\deleted{\unsure{of the terms of}}
% L02
i.e. the topos $B_{D,D^{*}}$ is equivalent to
% L03
$D_{\mathrm{ét}}$ -- now the $H^i_S(X_{\mathrm{ét}},F)$ are indeed calculated
% L04
on $D_{\mathrm{ét}}$\ldots We thus find

% L05
\uline{Proposition.} The homomorphism of exact cohomology
% L06
sequences under consideration is an isomorphism,
% L07
provided we know that the homomorphism
% L08
$H^i(\mathcal E,F_{\mathcal E})\xrightarrow{\sim}H^i(U_{\mathrm{ét}},F_U)$ is an isomorphism
% L09
for $\forall i$ ($F_U$ an \unsure{étale|étale local} sheaf on $U_{\mathrm{ét}}$
% L10
defined by a group $F_{\mathcal E}$ on which $\mathcal E=\pi_1(U)$
% L11
operates).

% L12
\uline{Example} : \inserted{OK} If $U$ is an affine curve
% L13
over algebraically closed $k$ \inserted{$F$ \deleted{\illegible{1}} prime to characteristic} (since OK for $i=0,1$, and
% L14
for $i=2$, the two \unsure{sides} are zero
% L15
(the Galois \unsure{cohomology|cohomology}, since the fundamental group
% L16
\inserted{prime to $p$} is \deleted{pro} free\ldots) -- whence we deduce the
% L17
analogous case for $\unsure{U|X}$ affine over arbitrary $k$,
% L18
then even if $U$ is not affine but
% L19
simply quasi-projective ? \deleted{\illegible{1}}
\margin{\uline{Corollary.} In this case, the cohomology of $X$ with coefficients in a locally constant $F_X$ \unsure{is} also calculated Galois-theoretically ? No, because of genus $0$ !!! In the case of smooth projective curves of genus $\ne0$ a special argument is needed.}

% L20
We are now interested in the
% L21
case where $X$ projective \inserted{(connected)} over algebraically closed $k$, \unsure{we thus see|we thus know}

\scanpage{130}
% L01
that : \deleted{the $H^i_!(U,F$}
% L02
a) If $S=\varnothing$ i.e. $X=U$, the cohomology of $U$
% L03
with (coefficients in local systems) is that
% L04
of $\pi$. Thus\ldots a canonical isomorphism
% L05
\[
H^2(\pi,T)\simeq H^2(X,T)\quad(\simeq\widehat{\mathbb Z})
\]
\margin{with \unsure{grain|grains} of salt in characteristic $p>0$}
% L06
\deleted{b) If $S=$}

% L07
b) If $S\ne\varnothing$, the $H^i_!(U,-)$ are \unsure{determined|described}
% L08
and calculated by Galois-theoretic means, in terms
% L09
of the \deleted{\unsure{systems}} \inserted{groups} with loops, making it possible to reconstruct
% L10
the situation
% L11
\diagram{\begin{tikzcd}[ampersand replacement=\&,column sep=large]
 \& B_{\pi_1,D^{*}} \arrow[dl] \arrow[dr] \& \\
 B_{\pi_1,U} \arrow[dr,dotted,no head] \& \&
 {\begin{gathered}
 B_{\pi_1,D}\ (={\text{discrete topos defined}}\\
 {\text{by }}S=I)
 \end{gathered}} \arrow[dl,dashed] \\
 \& B_{U,X} \&
\end{tikzcd}}
% L12
by postulating that the component $B_{\pi_1,D_i^{*}}$
% L13
is none other than $B_{L_i}$ (with the
% L14
\unsure{notations} of the preceding \unsure{number$^{o}$|case}). We
% L15
then find canonical isomorphisms (by local calculations)
% L16
\[
H^2_{\unsure{S_i|s_i}}(B_{U,X},\unsure{L_i})\simeq\widehat{\mathbb Z}
\quad\text{i.e.}\quad
H^2_{\unsure{S_i|s_i}}(B_{U,X},\widehat{\mathbb Z})\simeq L_i^{\otimes-1}
\]
% L17
(if $g\ne0$)
\[
H^2_{\unsure{S_i|s_i}}(B_{U,X},\widehat{\mathbb Z})
\xrightarrow{\sim}H^2(X,\widehat{\mathbb Z})
\simeq H^2(\pi,\widehat{\mathbb Z})
\]
% L18
whence \unsure{in} \deleted{\unsure{summary}} putting together, canonical
% L19
isomorphisms
\[
\deleted{\unsure{H}}\quad L_i\simeq H^2(\pi,\widehat{\mathbb Z})^{\otimes-1}.
\]
\margin{But in all cases (even if $g=0$, so long \unsure{as we} do not have $\nu=0$) \deleted{\illegible{1}} we find $H^{2}_{!}(B_{\pi_1,U},\widehat{\mathbb Z})\simeq\unsure{\widehat T_{\pi}^{\otimes-1}|T_{\pi}^{\otimes-1}}$ canonically, whence a \unsure{cohomological} description of the \unsure{orientation modules}, common to the \unsure{case without holes and with} holes\ldots}
